\chapter{Génération monodromique par survie de
\texorpdfstring{\(\pi\), \(e\) et \(\varphi\)}{pi, e et phi}}
\label{chap:generacion-numeros-rev6}
\label{chap:biografias-numeros-rev6}
\index{nombre HMT!génération par survie}
\index{pi@\(\pi\)!clôture régionale}
\index{e@\(e\)!propagation et mémoire}
\index{phi@\(\varphi\)!auto-échelle}
\index{retour nonadique}\index{miroir pi/e@miroir \(\pi/e\)}

La génération ne commence pas par un développement décimal prescrit. Elle commence dans
l'état complet du TPK : mot visible, feuillet, orientation, quotient,
retenue, état de Hensel, cylindres ternaire et décimal, signature conjointe,
frontière, phase et registre dodécaphasique. La relation EF--\(G_9\) prolonge cet
état et n'élimine une continuation que lorsqu'elle cesse de survivre à une
frontière future. Le lecteur décimal agit ensuite sur la section qui a
été déterminée.

La chaîne causale est
\[
\Omega_{\rm TPK}
\longrightarrow 468\,\mathcal U_6
\longrightarrow 243\,\mathcal W_6
\longrightarrow
\bigl(w_\pi,w_e,w_\varphi\bigr)
\longrightarrow
\bigl(\mathcal R_{\pi,n},\mathcal R_{e,n},
      \mathcal R_{\varphi,n}\bigr)_{n\geq0}
\longrightarrow
\varprojlim_n\mathcal R_{\chi,n}
\xrightarrow{\ L_{\rm arch}\ }
\chi .
\]
L'espace ambiant \(\mathbb F_3^6\) possède \(729\) mots, mais l'image
effective en contient \(243\) ; les deux objets ne sont pas identifiés. Au sein de cette
image, la sélection interne fixe
\[
w_\pi=010211,\qquad
w_e=201101,\qquad
w_\varphi=121200
\]
avant de consulter les chiffres de \(\pi,e,\varphi\).

\section{Régions génératives et section globale}

\begin{definition}[Région générative d'un canal]
\label{def:region-generativa-rev6}
Pour \(\chi\in\{\pi,e,\varphi\}\), soit \(\mathcal R_{\chi,n}\) l'ensemble
fini des états complets de profondeur \(n\) qui partent de
\(w_\chi\) et satisfont à chaque transition la relation EF--\(G_9\),
l'intersection exacte des cylindres et la signature conjointe de frontière. Si
\(\tau_{n+1,n}\) élimine uniquement la dernière couche, nous définissons
\[
\mathcal R_\chi
=\varprojlim_n
\bigl(\mathcal R_{\chi,n},\tau_{n+1,n}\bigr).
\]
La valeur publiée est \(L_{\rm arch}(x_\chi)\), où
\(x_\chi\in\mathcal R_\chi\) ; le lecteur ne fait pas partie de l'état qu'il
sélectionne.
\end{definition}

\begin{theorem}[Principe de prolongation monodromique dans le TPK complet]
\label{thm:generacion-supervivencia-pi-e-phi-rev6}
Supposons que, pour un canal
\(\chi\in\{\pi,e,\varphi\}\), chaque ensemble
\(\operatorname{Surv}_{\chi,n}\) soit non vide et unitaire, et que ses éléments
soient compatibles avec la troncature. Il existe alors une unique section
\[
x_\chi=(x_{\chi,n})_{n\geq0}\in\mathcal R_\chi
\]
Les cylindres décimaux associés à ses troncatures sont emboîtés et leurs
diamètres tendent vers zéro. Si le lecteur du canal identifie leurs intervalles
à la clôture, à la propagation ou à l'auto-échelle déclarées, l'intersection
publie, respectivement, \(\pi\), \(e\) ou \(\varphi\).
\end{theorem}

\begin{proof}
Dans le profil complet, la survie stable de niveau \(n\) est
\[
\operatorname{Surv}_{\chi,n}
=\bigcap_{h\geq1}
\left\{s:\text{\(s\) admet un prolongement compatible sur \(h\) frontières}\right\}.
\]
Par hypothèse, l'état survivant de chaque niveau se tronque en celui du
niveau précédent. Les niveaux forment un arbre à ramification finie et non vide
à toute profondeur. Le lemme de König fournit une branche globale ;
l'unitarité niveau par niveau empêche une seconde section.

Pour un préfixe ternaire de longueur \(L\), interprété comme \(N\), et un
préfixe de \(m\) triades, interprété comme \(D\), les compatibilités sont
\[
I_{N,L}=\left[\frac{N}{3^L},\frac{N+1}{3^L}\right),
\qquad
J_{D,m}=\left[\frac{D}{1000^m},\frac{D+1}{1000^m}\right),
\qquad
I_{N,L}\cap J_{D,m}\ne\varnothing .
\]
La prolongation conserve l'intersection et fait tendre vers zéro les deux
diamètres. L'intersection totale contient un unique réel. Le lecteur déclaré
identifie ce réel au sein de son canal.
\end{proof}

L'exécution reproductible ne remplace pas les hypothèses du théorème par une
affirmation implicite. Elle les vérifie à chacun des (3501) niveaux
nécessaires pour imprimer dix mille chiffres : la partition \(729\to1\), la
compatibilité des troncatures et le changement de mémoire après le retour y sont
contrôlés de manière exhaustive.

\section{Filtration de survie, monodromie nonadique et symétrie
\texorpdfstring{\(\pi/e\)}{pi/e}}

Pour
\[
B_n=(u_n^\pi,u_n^e,u_n^\varphi),\qquad
q_n=u_n^\pi+u_n^e+u_n^\varphi,\qquad
a_n=u_n^\pi+u_n^e-u_n^\varphi,
\]
la signature conjointe détermine
\[
u_n^\varphi=2(q_n-a_n),\qquad
u_n^\pi+u_n^e=q_n-u_n^\varphi
\quad\text{dans }\mathbb F_3^6.
\]
Ainsi, \(\varphi\) est l'axe fixe d'auto-échelle et l'unique ambiguïté locale
est la répartition spéculaire entre \(\pi\) et \(e\).

La phase est \(g(n)=n\bmod9\), avec \(0\equiv9\). Après \(w_{30}\), le
premier tour parcourt
\[
5,6,7,8,9,1,2,3,4.
\]
Ce ne sont pas neuf filtres indépendants, mais les neuf phases d'une même
connexion : élagage, synchronisation, orientation spéculaire, réouverture et
confirmation par survie.

À la neuvième phase,
\[
B_9=(100100\mid020112\mid010122)
\]
possède trois sorties localement compatibles. Les trois conservent
\(u_{10}^\varphi=112002\) et
\(u_{10}^\pi+u_{10}^e=210110\) ; toutes survivent à une frontière, mais seule
\[
\boxed{B_{10}=(101222\mid112221\mid112002)}
\]
survit à la seconde. Le futur brise le miroir local. De plus,
\[
g(n+9)=g(n),\qquad x_{\chi,n+9}\ne x_{\chi,n}
\]
lors des retours : l'étape suivante revient à la phase \(1\), non à l'état
précédent. Le quotient, la retenue et la frontière transforment le tour en une
hélice de mémoire. Cette action de retour avec transport non trivial est la
monodromie nonadique du TPK.

\section{Une émission reproductible de douze triades}

Treize blocs, équivalant à \(78\) trits, suffisent pour que le lecteur
\[
D=\left\lfloor\frac{1000^{12}N}{3^{78}}\right\rfloor
\]
publie la séquence suivante directement à partir des branches ternaires :
\[
\begin{array}{c|rrrrrrrrrrrr}
\pi&141&592&653&589&793&238&462&643&383&279&502&884\\
e&718&281&828&459&045&235&360&287&471&352&662&497\\
\varphi&618&033&988&749&894&848&204&586&834&365&638&117.
\end{array}
\]
Cette table n'alimente pas la transition : elle en montre la sortie. L'itération de la
même connexion, avec le retour \(9\to1\), produit tout préfixe fini
demandé sous les hypothèses de
\cref{thm:generacion-supervivencia-pi-e-phi-rev6} ; ce théorème précise la
prolongation à la limite inverse.

\section{Clôture unique de \texorpdfstring{\(\pi\)}{pi} et classification de ses \(5\) histoires régionales}

Soit
\[
 \Psi:\mathcal U_6\longrightarrow\mathbb F_3^6,
 \qquad
 \Psi(U_6)=(-U_6)\bmod3
\]
la projection coordonnée du catalogue. La fibre
\[
 \mathscr R_\pi^{\rm micro}
 :=\Psi^{-1}(010211)
\]
contient exactement les cinq vecteurs
\[
\begin{aligned}
U_\pi^\perp&=(501,614,498,169,272,272),\\
U_{\pi,1}^{++}&=(810,923,870,169,272,272),\\
U_{\pi,2}^{++}&=(810,983,810,169,272,272),\\
U_{\pi,3}^{++}&=(870,923,810,169,272,272),\\
U_\pi^{--}&=(870,923,810,418,674,521).
\end{aligned}
\]
Leurs signatures multiplicatives sont
\[
555555,\qquad339555,\qquad393555,\qquad933555,\qquad933717,
\]
et leurs multiplicités sont
\[
432,\qquad144,\qquad144,\qquad144,\qquad144.
\]
Par conséquent,
\[
\boxed{
5=1_\perp+3_{++}+1_{--},
\qquad
1008=432+4\cdot144.}
\]
La première identité classe des fonctions d'orientation ; la seconde compte
des routes germes. Ce ne sont pas deux manières de compter cinq copies du même objet.

Le quotient visible
\[
 \mathscr R_\pi^{w_6}
 =\Psi(\mathscr R_\pi^{\rm micro})
 =\{010211\}
\]
possède un seul élément car il oublie la provenance régionale. L'égalité de
mot n'identifie pas les cinq états : les quatre premières régions
conservent la queue \(169|272|272\), tandis que la région négative réalise un
retour muté. La région transversale, de multiplicité \(432\), ne se
réduit pas non plus à l'une des quatre branches de multiplicité \(144\).

La fonction de fermeture se voit également dans le multigraphe des blocs. Chacune des
cinq régions appartient à une marche fermée minimale de longueur huit qui
parcourt aussi les ancrages \(U_e\), \(U_\varphi\) et \(U_{\rm APP}\).
L'énumération et la démonstration de minimalité se trouvent dans
\cref{chap:cinco-ciclos-pi} ; la liste régionale et ses multiplicités, dans
\cref{chap:regiones}.

\begin{proposition}[Paquet de clôture de \(\pi\)]
\label{prop:paquete-clausura-pi-rev6}
La donnée
\[
 \mathfrak C_\pi=
 \bigl(
 \mathscr R_\pi^{\rm micro},
 w_\pi,
 \rho_\pi,
 m_\pi,
 \ell_{\rm cl}
 \bigr),
\]
où
\[
\rho_\pi=(\perp,++ ,++ ,++ ,--),\qquad
m_\pi=(432,144,144,144,144),\qquad
\ell_{\rm cl}=8,
\]
conserve exactement l'information fonctionnelle que perd le quotient
\(\mathscr R_\pi^{w_6}\). Par conséquent, la fonction de \(\pi\) dans cette
génération est de clore cinq histoires régionales dans une seule évaluation,
non de répéter cinq fois une même constante.
\end{proposition}

\begin{proof}
La fibre, les rôles et les multiplicités sont ceux de
\cref{chap:regiones}. \cref{chap:cinco-ciclos-pi} démontre que chaque
région entre dans une fermeture minimale de longueur huit avec les trois ancrages
déclarés. La projection \(\Psi\) envoie les cinq régions sur le même
mot, de sorte que le quotient conserve la valeur visible et oublie
\(\rho_\pi,m_\pi,\ell_{\rm cl}\).
\end{proof}

\section{\texorpdfstring{\(e\)}{e} : propagation visible et mémoire de quotient différenciée}

La fibre de \(w_e=201101\) se compose d'une émission \(U_e\) de multiplicité
\(144\). Son trait fonctionnel n'est pas une périodicité infinie, mais une
répétition locale distinguée exhaustivement au sein du lecteur des
\(243\) mots. \cref{thm:cuadrado-local-e} démontre que \(w_e\) est
l'unique mot dont la lecture de douze chiffres contient un carré de longueur
quatre commençant au deuxième chiffre :
\[
718281828459
=7\,\underbrace{1828\,1828}_{\text{carré local}}\,459.
\]
La juxtaposition de blocs est une concaténation décimale ; ce n'est pas le produit
arithmétique \(7\cdot1828^2\cdot459\). Le catalogue contient un autre carré de
longueur quatre,
\[
575157518472
=\underbrace{5751\,5751}_{\text{carré local}}\,8472,
\]
mais il commence au premier chiffre et ne satisfait donc pas le même
prédicat positionnel.

Le recensement complet des carrés de longueurs \(1,\ldots,6\) est
\[
\begin{array}{crrrrrr}
\ell&1&2&3&4&5&6\\
\text{occurrences}&240&21&3&2&0&0\\
\text{mots qui les contiennent}&153&19&3&2&0&0.
\end{array}
\]
La répétition \(1828\,1828\) n'amorce pas non plus une période : une troisième copie
exigerait comme chiffre suivant un \(1\), mais la lecture effective contient
un \(4\).

La dynamique d'élévation montre pourquoi l'épisode exprime bien une
propagation :
\[
201101\mid121221\mid102011\mid012222\mid102011.
\]
Les troisième et cinquième blocs visibles coïncident, mais les états qui les
émettent sont distincts. Leurs pré-états modulo \(27\) sont
\[
(25,11,7,17,8,16),
\qquad
(7,8,4,23,26,7),
\]
et leurs quotients de Hensel ultérieurs sont
\[
(6,13,9,2,13,1),
\qquad
(15,0,3,19,23,25).
\]
Il y a retour de la projection \(102011\), non retour de l'état enrichi.
La sortie peut se propager et réapparaître parce que le quotient transporte la
généalogie que le mot visible a omise.

\begin{proposition}[Contenu fonctionnel de \(e\)]
\label{prop:propagacion-e-rev6}
La génération exacte de \(e\) réunit trois faits compatibles :
\begin{enumerate}
\item la sélection interne de \(w_e=201101\) ;
\item l'unicité positionnelle du carré local \(1828\,1828\) dans le
catalogue fixé ;
\item la réapparition de \(102011\) sous des états de mémoire différents.
\end{enumerate}
Sa fonction est la propagation avec renouvellement de mémoire. Ce n'est ni la périodicité du
développement ni la répétition de l'état complet.
\end{proposition}

\begin{proof}
Le premier point est le cas \(e\) de
\cref{thm:selector-orbital-constantes} ; le second et son recensement sont
\cref{thm:cuadrado-local-e}. L'inégalité des deux pré-états et des
deux quotients de Hensel démontre le troisième.
\end{proof}

\section{Axe d'auto-échelle et fibre spéculaire}

La signature conjointe des trois générations conserve une décomposition linéaire
qui sera utilisée ensuite par le lecteur d'échelle. Pour
\[
 B_k=(u_k^\pi,u_k^e,u_k^\varphi)
 \in(\mathbb F_3^6)^3
\]
nous définissons, coordonnée par coordonnée,
\[
 q_k=u_k^\pi+u_k^e+u_k^\varphi,
 \qquad
 a_k=u_k^\pi+u_k^e-u_k^\varphi.
\]

\begin{theorem}[Décomposition axiale de la signature de frontière]
\label{thm:eje-espejo-nonadico}
La paire \((q_k,a_k)\) détermine univoquement le canal d'auto-échelle et
l'agrégat des deux autres canaux :
\[
 u_k^\varphi=2(q_k-a_k),
 \qquad
 u_k^\pi+u_k^e=q_k-u_k^\varphi.
\]
Par conséquent, une fibre de
\[
 q_{\rm ax}(u^\pi,u^e,u^\varphi)
 =(u^\pi+u^e,u^\varphi)
\]
ne conserve d'ambiguïté que dans la répartition interne entre \(u^\pi\) et \(u^e\).
\end{theorem}

\begin{proof}
Soustraire les deux équations de définition donne
\(q_k-a_k=2u_k^\varphi\). Comme l'inverse de \(2\) dans \(\mathbb F_3\) est
à nouveau \(2\),
\[
u_k^\varphi=2(q_k-a_k).
\]
En soustrayant ce vecteur de \(q_k\), il reste
\(u_k^\pi+u_k^e\). L'opération fixe l'axe \(\varphi\), mais ne sépare pas
les deux termes de la fibre \(\pi/e\) ; c'est exactement l'ambiguïté
résiduelle.
\end{proof}

\def\HMTAxialTheoremDefined{1}

L'expression \emph{fibre spéculaire} désigne cette décomposition exacte dans
\(\mathbb F_3^6\). Elle ne postule pas une identité analytique entre les constantes
réelles : elle indique quelle composante est fixée par la signature et où subsiste un
choix interne.

\section{\texorpdfstring{\(\varphi\)}{Phi} : l'auto-échelle induite par TPK}

Le mot \(w_\varphi=121200\) possède une seule émission de multiplicité
\(144\), mais sa fonction ne se déduit pas de ce cardinal. Le calendrier TPK
produit le bloc de modes
\[
(\Sigma,\Pi,\Pi).
\]
En lisant \(\Sigma\) comme feuillet aréal et \(\Pi\) comme feuillet volumétrique, les
trois paires consécutives du cycle imposent l'incidence
\[
F_{\rm av}
=
\begin{pmatrix}
0&1\\
1&1
\end{pmatrix},
\]
comme le démontre \cref{thm:descenso-necesario-fav}. Le nombre d'or n'est pas
introduit pour choisir cette matrice : il apparaît ensuite comme sa racine de Perron.
En effet,
\[
\det(tI-F_{\rm av})=t^2-t-1,
\qquad
F_{\rm av}\binom{1}{\varphi}
=\varphi\binom{1}{\varphi}.
\]

Pour un chemin admissible \(\gamma:i\to j\) de longueur \(n\), avec
\(r=(1,\varphi)\), le caractère de Parry
\[
\chi_P(\gamma)=\varphi^n\frac{r_i}{r_j}
\]
est multiplicatif par concaténation. Sur un chemin fermé,
\(\chi_P(\gamma)=\varphi^n\). La mesure d'entropie maximale de l'enveloppe
à mémoire un satisfait, d'après \cref{cor:parry-envolvente-fav},
\[
\frac{\mu_{\rm ar}}{\mu_{\rm vol}}=\varphi^{-2}.
\]

La mémoire du second ordre conserve en outre la direction contractée. Dans la
base \((x^2,2xy,y^2)\),
\[
\operatorname{Sym}^2(F_{\rm av})
=
\begin{pmatrix}
0&0&1\\
0&1&2\\
1&1&1
\end{pmatrix},
\qquad
\operatorname{Spec}\bigl(\operatorname{Sym}^2(F_{\rm av})\bigr)
=\{\varphi^2,-1,\varphi^{-2}\}.
\]
La valeur propre \(\varphi^{-2}\) est le mode stable positif. Sur des retours
fermés, le cocycle d'échelle de \cref{sec:cociclo-parry-escala} donne
\[
M_n=\varphi^{-2n}={\bigl(\chi_P(\gamma)\bigr)}^{-2}.
\]

\begin{proposition}[Contenu fonctionnel de \(\varphi\)]
\label{prop:autoescala-phi-rev6}
La fonction de \(\varphi\) dans HMT est l'auto-échelle du quotient de modes :
croissance par la racine de Perron \(\varphi\), conservation cylindrique par
\(\chi_P\) et retour stable du second ordre par \(\varphi^{-2}\).
Ces trois expressions appartiennent à un même opérateur induit par le
calendrier TPK.
\end{proposition}

\begin{proof}
L'incidence \(F_{\rm av}\) descend du bloc
\((\Sigma,\Pi,\Pi)\) par
\cref{thm:descenso-necesario-fav}. Son polynôme caractéristique produit la
racine de Perron et la multiplicativité de \(\chi_P\) s'obtient par
télescopie. L'action quadratique sépare les trois valeurs propres indiquées ; le
mode positif contracté est \(\varphi^{-2}\).
\end{proof}

\section{Correspondance orientée entre clôture et auto-échelle}

La matrice \(F_{\rm av}\) engendre \(\varphi^{-2}\), mais ne peut pas engendrer à
elle seule \(\pi\) : toute fonction rationnelle de son spectre appartient à
\(\mathbb Q(\sqrt5)\). La clôture transcendante provient du lecteur régional
coinductif. Cette séparation des origines est conservée par l'opérateur
diagonal
\[
 D_\pi=\pi_{\rm HMT}P_{\rm ar}+P_{\rm vol},
\qquad
 P_{\rm ar}=
 \begin{pmatrix}1&0\\0&0\end{pmatrix},
\quad
 P_{\rm vol}=
 \begin{pmatrix}0&0\\0&1\end{pmatrix}.
\]
Le retour marqué de longueur \(n\) est
\[
\Theta_n
=\pi_{\rm HMT}\frac{F_{n-1}}{F_{n+1}},
\qquad
\lim_{n\to\infty}\Theta_n
=\frac{\pi_{\rm HMT}}{\varphi^2}.
\]
La démonstration dynamique, le témoin exact de longueur \(108\) et le lecteur
de cylindres emboîtés se trouvent dans
\cref{sec:retorno-marcado-pi-phi2}. La marque de \(\pi\) est appliquée une seule
fois au retour : elle n'est pas insérée dans chaque transition de \(F_{\rm av}\).

Le même rapport possède deux orientations. Soient
\[
\mathscr R(x,y)=\frac{x}{y^2},
\qquad
\mathsf J(x,y)=(y,x).
\]
Alors \(\mathsf J^2=I\) et
\[
\boxed{
\mathscr R(\pi,\varphi)=\frac{\pi}{\varphi^2},
\qquad
(\mathscr R\circ\mathsf J)(\pi,\varphi)
=\frac{\varphi}{\pi^2}.}
\]
De plus,
\[
\frac{\pi/\varphi^2}{\varphi/\pi^2}
={\left(\frac{\pi}{\varphi}\right)}^3,
\qquad
\frac{\varphi}{\pi^2}
=\frac{1}{\varphi^3{\bigl(\pi/\varphi^2\bigr)}^2}.
\]
La première orientation lit la clôture aréale sur la mémoire volumétrique ; la
seconde, l'auto-échelle intérieure sur la fermeture quadratique. La
section~\ref{sec:espejo-pi-phi-electron-hmt} démontre cette involution et son
apparition ultérieure dans l'exposant électronique. Cette correspondance n'identifie pas
\(\pi\) à \(\varphi\) : elle conserve leur ordre et explicite ce qui change en
l'inversant.

\section{Précision arbitraire de l'émission monodromique}
\label{sec:precision-arbitraria-monodromia-rev6}

\begin{corollary}[Émission de tout préfixe fini]
\label{cor:precision-arbitraria-monodromia-rev6}
Sous les hypothèses de
\cref{thm:generacion-supervivencia-pi-e-phi-rev6}, pour tout \(M\geq1\),
l'itération de EF--\(G_9\) détermine les premières
\(M\) décimales de chacun des trois canaux. En particulier, \(M=100\)
et \(M=1000\) sont des fenêtres finies de la même prolongation monodromique et ne
nécessitent pas l'introduction d'une table extérieure de chiffres.
\end{corollary}

\begin{proof}
Soit \(K(n)\) le nombre de triades publiées après \(n\) blocs
ternaires. La section fournie par le théorème possède des niveaux de toute
profondeur et \(K(n)\to\infty\). Étant donné \(M\),
choisissons \(n\) avec \(3K(n)\geq M\).
\cref{thm:generacion-supervivencia-pi-e-phi-rev6} détermine de manière unique
la troncature de profondeur \(n\), et son cylindre décimal fixe les
\(3K(n)\) premiers chiffres. Prendre les \(M\) premiers conclut.
\end{proof}

La différence des fonctions est fixée au sein d'une seule construction :
\(\pi\) clôt des provenances régionales distinctes dans une évaluation commune ;
\(e\) propage une sortie visible tout en renouvelant la mémoire qui la
transporte ; \(\varphi\) régit l'auto-échelle du quotient de modes. La
monodromie explique comment ces fonctions reviennent à la même phase et traversent
le miroir \(\pi/e\). Sous les hypothèses coinductives du théorème de
prolongation, elles publient indéfiniment de nouveaux blocs sans réinitialiser leur
histoire.

% El resultado de realizabilidad de una cinta prescrita se conserva como
% material auxiliar, pero no forma parte del argumento generativo rector.
\ifdefined\HMTIncludeTargetFedRay
\section{Prolongation coinductive par rayons : existence, compatibilité et domaine}
\label{sec:rayo-coinductivo-rev6}

Le passage de tout préfixe fini à une route infinie utilise un théorème de
compacité à ramification finie. Son objet n'est pas le sélecteur du catalogue,
mais la réalisation d'un canal dont le développement a déjà été fixé.

Soit
\[
\mathcal S
=\mathbb Z_9^2
\times\{+,\times\}
\times\{N,E,S,O\},
\qquad
|\mathcal S|=648,
\]
l'espace fini d'états du modèle de routes. L'alphabet de contrôle est
\(\mathcal A=\{N,E,S,O\}\), et le coût d'une route compte les rotations :
\[
c(a_{t-1},a_t)=\mathbf1_{\{a_t\neq a_{t-1}\}}.
\]
Une fois fixés un canal \(x\in\{\pi,e,\varphi\}\) et son développement ternaire, soit
\(W_m(x)\) une route canonique d'horizon \(27m\) qui minimise ce coût
parmi les routes qui réalisent les \(9m\) premiers trits de \(x\).

\begin{theorem}[Rayon limite d'un développement fixé]
\label{thm:rayo-coinductivo-condicionado-rev6}
Pour chaque développement fixé \(x\in\{\pi,e,\varphi\}\), il existe une route
\[
\mathcal D_\infty(x)\in\mathcal A^{\mathbb N}
\]
tel que :
\begin{enumerate}
\item pour tout \(L\geq1\), le préfixe de longueur \(L\) de
\(\mathcal D_\infty(x)\) coïncide avec le préfixe de \(W_m(x)\) pour
une infinité de valeurs de \(m\) ;
\item \(\mathcal D_\infty(x)\) réalise exactement le développement ternaire
infini de \(x\).
\end{enumerate}
Parmi les branches ayant la première propriété, on prend la branche lexicographiquement
minimale.
\end{theorem}

\begin{proof}
Pour chaque \(L\), considérons l'ensemble \(T_L(x)\) des préfixes de longueur
\(L\) qui apparaissent au début de \(W_m(x)\) pour une infinité de valeurs de
\(m\). Il est non vide : il n'y a qu'un nombre fini de préfixes de longueur \(L\) et il y a
une infinité d'horizons. Les ensembles \(T_L(x)\), ordonnés par extension,
forment un arbre à ramification finie. Le lemme de König produit une branche
infinie ; la plus petite dans l'ordre lexicographique définit
\(\mathcal D_\infty(x)\).

Une fois \(n\) fixé, prenons \(L\) suffisamment grand pour déterminer les \(n\) premiers
trits de filtration. Ce préfixe apparaît dans une infinité de \(W_m(x)\), et chacun
d'entre eux réalise par définition les \(n\) premiers trits de \(x\). Le même
préfixe de \(\mathcal D_\infty(x)\) les réalise. Comme \(n\) est arbitraire,
la route réalise le développement complet.
\end{proof}

Soit maintenant
\[
\mathcal R(x)=\{729x\},
\]
où \(\{\cdot\}\) désigne la partie fractionnaire. Multiplier par
\(729=3^6\) décale six trits. Dans le calendrier du modèle, chaque trit de
filtration occupe trois itérations élémentaires, de sorte que six trits
équivalent à dix-huit itérations élémentaires.

\begin{corollary}[Covariance par l'auto-application \(729\)]
\label{cor:covarianza-rayo-729-rev6}
Soit \(\sigma\) le décalage vers la gauche d'une route. Pour tout
\(r\geq0\),
\[
\sigma^{18r}\mathcal D_\infty(x)
\]
réalise exactement le canal \(\mathcal R^r(x)\).
\end{corollary}

\begin{proof}
La covariance finie identifie le suffixe qui commence après \(18r\)
itérations élémentaires
au décalage de \(6r\) trits.
\cref{thm:rayo-coinductivo-condicionado-rev6} transporte cette identité à
tous les préfixes du rayon limite.
\end{proof}

Le corollaire exprime une holographie itérée : une même cellule finie,
décalée de \(\mathcal R\), réalise des strates successives du développement.
Il ne transforme pas \(\mathcal R\) en sélecteur de \(\pi,e,\varphi\), car
\(x\) intervient déjà dans la définition des ensembles de routes admissibles
\(W_m(x)\). Sa fonction exacte est ultérieure : une fois le canal distingué,
il construit un rayon compatible et rend visible sa covariance entre échelles.

\section{Composition du sélecteur orbital, du cylindre numérique et de l'évaluation archimédienne}

Les trois opérations finales forment un diagramme causal, non trois versions
du même théorème :
\[
\text{état interne}
\xrightarrow{\,\mathfrak S_{\rm cat}\,}
w_x
\xrightarrow{\ \text{lecteur enrichi}\ }
\text{développement }x
\xrightarrow{\,\mathcal D_\infty(x)\,}
\text{rayon qui la réalise}.
\]
Le sélecteur orbital \(\mathfrak S_{\rm cat}\) distingue les trois mots
par des invariants du catalogue et un calibre d'orientation explicite. Le lecteur
enrichi prolonge le mot, le cylindre et la mémoire. Le théorème des rayons
construit ensuite une réalisation infinie optimale récurrente du développement
déjà fixé.

\begin{proposition}[Séparation de force entre les trois couches]
\label{prop:tipado-biografias-rayos-selector-rev6}
Les affirmations suivantes sont simultanément satisfaites.
\begin{enumerate}
\item Les cinq régions et les fermetures minimales de \(\pi\), le carré local
et la rupture de mémoire de \(e\), et l'auto-échelle
\(F_{\rm av}\)--\(\varphi\)--\(\varphi^{-2}\) sont des résultats exacts dans
leurs domaines déclarés.
\item Le rayon \(\mathcal D_\infty(x)\) est une réalisation coinductive exacte
d'un développement fixé.
\item Le rayon n'est pas un sélecteur autonome, car le développement \(x\) apparaît
dans la définition de \(W_m(x)\).
\end{enumerate}
La troisième affirmation limite uniquement l'usage sélecteur du théorème des
rayons ; elle ne modifie pas les deux premières.
\end{proposition}

\begin{proof}
Les résultats du premier point sont démontrés, respectivement, dans
\cref{chap:regiones,chap:cinco-ciclos-pi,thm:cuadrado-local-e,%
thm:descenso-necesario-fav,cor:parry-envolvente-fav}. Le second point est
\cref{thm:rayo-coinductivo-condicionado-rev6}. Pour le troisième, il suffit
d'observer que l'ensemble de minimisation qui définit \(W_m(x)\) exige de
réaliser les \(9m\) premiers trits de \(x\) ; on ne peut donc pas utiliser
cette même route pour déduire quel \(x\) devait être sélectionné.
\end{proof}

La différence des fonctions est finalement fixée. \(\pi\) clôt des
histoires régionales distinctes dans une seule évaluation ; \(e\) propage une
sortie visible tout en renouvelant la mémoire qui la transporte ; \(\varphi\)
régit le changement d'échelle du quotient de modes. La correspondance
\(\pi/\varphi^2\leftrightarrow\varphi/\pi^2\) coordonne fermeture et
autocroissance sans les confondre. Le rayon coinductif explique comment un
développement déjà distingué se prolonge indéfiniment. C'est cela, la biographie :
non seulement quel nombre apparaît, mais quelle structure le sélectionne, quelle mémoire
il conserve et quelle fonction il remplit dans le système.
\fi
